\section{Interlude: RV decomposition and Fubini theorems}\label{csndag}

We continue to work in the original language $\lan{}{RV}{}$, over the
standing parameter base $\mathbb S$. In particular, no cross-section is
added to the language. We use the decomposition below and the
retraction $\mathbb E_b$ of Section~4 to prove the Fubini statements.

For an integer matrix $M=(m_{ab})\in\GL_n(\Z)$, write
\[
 h_M:(\RV^\times)^n\longrightarrow(\RV^\times)^n,\qquad
 (h_M(u))_a=\prod_b u_b^{m_{ab}}.
\]
Thus $h_M^{-1}=h_{M^{-1}}$ and
$\vrv(h_M(u))=M\vrv(u)$. For $I\subseteq\Gamma^\ell$, the notation
$I^\sharp$ in the following lemma means its full inverse image in
$(\RV^\times)^\ell$ under $\vrv$. The singleton zero/infinity element
of $\RV_\infty$ is treated separately by fixing its coordinate positions.

\setcounter{thm}{11}
\begin{lem}[Monomial decomposition]\label{coh:decom}
Let $U\subseteq(\RV^\times)^n$ be definable. There is a finite definable
partition $(D_i)_i$ of $\vrv(U)$ such that, writing
$U_i=U\cap\vrv^{-1}(D_i)$, there are $M_i\in\GL_n(\Z)$,
$0\leq m_i\leq n$, a nonempty definable set
$V_i\subseteq(\RV^\times)^{m_i}$, and a nonempty definable set
$I_i\subseteq\Gamma^{n-m_i}$ with
\[
 \vrv(V_i)=\{\beta_i\},\qquad
 h_{M_i}(U_i)=V_i\times I_i^\sharp .
\]
All sets and maps are definable over the original base $\mathbb S$.
In particular, the partition of $U$ is induced by a partition of
its value support.
\end{lem}

\begin{proof}
We give the normal-form step as well as the integer reduction.

\emph{Finite angular data.}
For $e\in\Z^n$, put $u^e=\prod_j u_j^{e_j}$.
The guarded RV normal form obtained from quantifier elimination has
the following consequence. There is a finite definable partition
$(P_a)_a$ of $\Gamma^n$ and, for each $a$, a finite set
$E_a\subseteq\Z^n$ such that, for $X_a=U\cap\vrv^{-1}(P_a)$:
\begin{enumerate}
\item for every $e\in E_a$, the set
$\{\vrv(u^e):u\in X_a\}$ is finite;
\item if $u\in X_a$, $\vrv(u)=\vrv(u')$, and
$u^e=(u')^e$ for every $e\in E_a$, then $u'\in X_a$.
\end{enumerate}
Here is the finite syntactic argument for this consequence.
In the guarded normal form, a monomial is $c u^e$ and a residue
polynomial is a finite sum of such monomials on their residue
domains. Partition the value space according to the residue-domain
status of every summand and the equality of the values of the two
outer monomials in every equality. Retain all nonzero and domain
guards and both signs of every literal; this partitions the entire
Boolean formula, rather than just covering it by true clauses.
A nonzero summand contributing to a residue polynomial forces
\[
 e\cdot\vrv(u)=-\vrv(c).
\]
Include its exponent $e$ among the active characters. A zero
coefficient contributes nothing and needs no character.
In an equality whose two nonzero outer coefficients have the same
value, include the difference of their exponent vectors: its value
is fixed by the quotient of those coefficients. If their values are
different, the equality is possible precisely when both residue
factors vanish. Thus every remaining angular test depends only on
these finitely many characters. Order tests depend only on the
coordinate values. This proves (1) and (2), including negated
literals.

The value partition is defined over $\mathbb S$. To check this if RV
representatives of named value-group parameters are used while
obtaining the normal form, the values of nonzero coefficients
descend by the original-base Gamma parameter-control theorem:
the value of a nonzero coefficient definable from these names
lies in $\dcl^{\mathrm{eq}}(\mathbb S)\cap\Gamma$.
The partition uses only those coefficient values. No chosen
RV representative is retained as a new parameter.

\emph{Integer coordinates.}
Fix $a$ and let $N$ have rows the characters in $E_a$.
Column Hermite reduction supplies $A\in\GL_n(\Z)$ with
\[
 NA=(C\ \ 0),\qquad \operatorname{rank}_{\mathbb Q}C=m,
\]
where $C$ has $m$ columns. This is an integer change of basis:
column additions, interchanges and sign changes suffice.
Set $z=h_{A^{-1}}(u)=(v,w)$ with $v\in(\RV^\times)^m$.
By (1), $C\vrv(v)$ has finite image on $h_{A^{-1}}(X_a)$.
The map $\beta\mapsto C\beta$ on $\Gamma^m$ is injective, since
$C$ has full column rank and the ordered group $\Gamma$ is
torsion-free. Consequently $\vrv(v)$ has finite image as well.
These finitely many tuples are individually $\mathbb S$-definable:
a finite definable subset of the ordered set $\Gamma^m$ can be
enumerated lexicographically. Refine $P_a$ by these value tuples
and fix one, denoted by $\beta$.

On this piece, membership of $(v,w)$ depends only on $v$ and
$\vrv(w)$. Indeed, changing $w$ within its value fiber keeps all
coordinate values of $u=h_A(v,w)$ unchanged, while $NA=(C\ 0)$
keeps every active character unchanged. Property (2) applies in
both directions. Therefore this piece is the full inverse image
of a definable relation
\[
 T\subseteq\beta^\sharp\times\Gamma^{n-m}
\]
under $(v,w)\mapsto(v,\vrv(w))$.

\emph{Separate the residue and value data over the original base.}
Temporarily choose a point of $\beta^\sharp$ to identify this
fiber with a residue-field torus. Residue-field/Gamma
orthogonality then implies that only finitely many distinct sets
\[
 T_\eta=\{v\in\beta^\sharp:(v,\eta)\in T\},
 \qquad \eta\in\Gamma^{n-m},
\]
occur. This finiteness does not retain the chosen point as a
parameter.

We spell out descent to $\mathbb S$. Put
$\Gamma_0=\dcl^{\mathrm{eq}}(\mathbb S)\cap\Gamma$.
The equivalence relation
\[
 \eta\sim\eta' \quad\Longleftrightarrow\quad
 \forall v\in\beta^\sharp\
 \big((v,\eta)\in T\ \longleftrightarrow\ (v,\eta')\in T\big)
\]
is $\mathbb S$-definable. By the original-base Gamma
parameter-control consequence of quantifier elimination, it is
definable in the pure ordered divisible group over $\Gamma_0$.
The standing nontriviality assumption makes $\Gamma_0$ a
nontrivial divisible ordered subgroup, so
$\Gamma_0\preceq\Gamma$. If there are $q$ equivalence classes,
the assertion that $q$ pairwise inequivalent representatives
exist is a first-order assertion over $\Gamma_0$. Hence we may
take representatives $\eta_1,\ldots,\eta_q\in\Gamma_0^{n-m}$.
Set
\[
 I_j=\{\eta:\eta\sim\eta_j\},\qquad V_j=T_{\eta_j}.
\]
These sets are definable over $\mathbb S$, the $I_j$ form a
partition, and
$T\cap(\beta^\sharp\times I_j)=V_j\times I_j$.
Discard the empty residue fibers. Each remaining piece has
full RV preimage $V_j\times I_j^\sharp$ and
$\vrv(V_j)=\{\beta\}$.

The corresponding original value piece is
$D_j=A(\{\beta\}\times I_j)$. It lies in the chosen $P_a$:
choose any $v\in V_j$ and lift each $\eta\in I_j$ to see this.
The $D_j$ partition the value support of the fixed-beta piece
of $X_a$.
Doing this for each $a$ proves the assertion with
$M_i=A^{-1}$. The cases $m=0$, $m=n$ and $U=\varnothing$
are covered by the usual empty-tuple conventions.
\end{proof}

% BEGIN REFEREE INSERTION R080
\ResolvedRefereeRemark{080}{Error}{Integer column operations}{%
\textit{Original source line 1789.}\quad The parenthetical assertion about $\mathbb Z$ is incorrect: every matrix in $\mathrm{GL}_n(\mathbb Z)$ is a product of elementary integer column additions, column interchanges, and sign changes. One may simply invoke column Hermite reduction to obtain $N_iM_i=(C_i\;0)$ with $C_i$ of full column rank. Then $C_i\vrv(v)=\alpha_i$ determines the first coordinates uniquely, which supplies the asserted product decomposition. Also, the cited item 3.25 of \cite{hrushovski:kazhdan:integration:vf} is a corollary, and ``Lemmas~3.21'' should be ``Lemma~3.21''.
}{The author approved the replacement on November 13, 2026. Lemma~\ref{coh:decom} now has an explicit original-base monomial decomposition and a proof from guarded RV normal form, integer reduction, and residue/Gamma separation. The incorrect parenthetical and the assumed HK 3.21/3.25 input are removed.}
% END REFEREE INSERTION R080
